\chapter{Central characters and blocks}
\label{chap:library:blocks}

A central idempotent gives a direct product decomposition of a ring.
For a group algebra, primitive central idempotents define the blocks.
Central subgroup idempotents relate the scalar action on a simple module
to its block. The Brauer homomorphism and block induction relate blocks
of different groups. They give Brauer's first main theorem under the
stated facts about defect groups. We also bound the group of linear
characters stabilising a block.

Throughout, $G$ is a finite group. Unless a stronger assumption is stated,
$k$ is a field. In arguments of positive characteristic, $\ell$ denotes a
prime and $\operatorname{char}k=\ell$. For
$z=\sum_{g\in G}z_g g\in kG$, the symbol $z_g$ denotes its coefficient at
$g$. A block means a primitive central idempotent of $kG$.
When blocks are indexed by a finite set $\mathcal B(G)$, their idempotents
are denoted by $b_B$, for $B\in\mathcal B(G)$. A complete primitive
decomposition and the central characters of its blocks will be assumed
where they are needed.

\section{Central idempotents and simple modules}
\label{lib:blocks:support}

The initial arguments apply to an arbitrary unital ring $R$.

\begin{definition}
A \emph{complete family of orthogonal central idempotents} is a finite family
$(e_i)_{i\in I}$ in $R$ such that each $e_i$ is central,
$e_i^2=e_i$, $e_i e_j=0$ for $i\ne j$, and $\sum_i e_i=1$.
A nonzero central idempotent $b$ is \emph{primitive} if every central
idempotent $c$ with $cb=c$ is either $0$ or $b$.
A module \emph{belongs to the block} $b$ if $b$ acts as the identity.
\end{definition}

\begin{proposition}[Central idempotents on simple modules]
\label{lib:blocks:unique-support}
Let $(e_i)_{i\in I}$ be a complete family of orthogonal central
idempotents in $R$.
Every simple $R$-module $V$ has a unique index $i$ such that $e_i$ acts
as the identity. Every other $e_j$ acts as zero.
Every primitive central idempotent $b$ likewise has a unique index $i$
such that
\[
 be_i=b,\qquad be_j=0\quad(j\ne i).
\]
If $b$ acts as the identity on $V$, the two indices coincide.
\end{proposition}
\begin{proof}
Centrality makes $v\mapsto e_i v$ an $R$-linear endomorphism. On a
simple module it is either zero or bijective. In the latter case its
idempotence implies that it is the identity: injectivity applied to
$e_i(e_i v)=e_i v$ gives $e_i v=v$. If every $e_i$ acted as zero, then
$1=\sum_i e_i$ would act as zero on the nonzero module $V$. Thus at least
one acts as the identity. If $e_i$ does, then
$e_jv=e_je_iv=0$ for $j\ne i$, proving uniqueness.

For the assertion about $b$, each $be_i$ is a central idempotent with
$(be_i)b=be_i$. Primitivity makes it either zero or $b$. The equation
$\sum_i be_i=b\ne0$ gives a nonzero term. If $be_i=b$, multiplication
by $e_j$ and orthogonality give $be_j=0$ for $j\ne i$.
Finally, if $b$ acts as the identity, then
$e_jv=b(e_jv)=(be_j)v$, so the two indices coincide.
\end{proof}

\begin{implementation}
The structure for a complete family of central idempotents extends Mathlib's
\lean{CompleteOrthogonalIdempotents}. The primitivity predicate is local,
and the proof for simple modules uses Mathlib's Schur lemma.
The two existence and uniqueness theorems are
\lean{ModularRep.CompleteOrthogonalCentralIdempotents.existsUnique_isSupport}
and \lean{ModularRep.IsPrimitiveCentralIdempotent.existsUnique_isSupport}.
The agreement of the two indices is
\lean{ModularRep.IsPrimitiveCentralIdempotent.IsSupport.toModuleSupport}.
Their modules are \module{ModularRep.CentralIdempotentSupport} and
\module{ModularRep.PrimitiveCentralIdempotent}.
\end{implementation}

A \emph{block decomposition} is a complete family of orthogonal
primitive central idempotents. For the expansion of central idempotents in
the group algebra case, see, for example,
\cite[Theorem~3.11]{Navarro1998}.

\begin{proposition}[Expansion in primitive idempotents]
\label{lib:blocks:idempotent-expansion}
Let $(b_B)_{B\in\mathcal B}$ be a block decomposition of $R$, and let
$e$ be a central idempotent. Set $S(e)=\{B\in\mathcal B:b_Be\ne0\}$. Then
\[
 b_Be\in\{0,b_B\},\qquad
 e=\sum_{B\in S(e)}b_B.
\]
The set $S(e)$ is the unique subset of $\mathcal B$ with this sum.
In particular, $e\ne0$ if and only if $S(e)$ is nonempty.
\end{proposition}
\begin{proof}
The product $b_Be$ is a central idempotent below $b_B$, so primitivity
gives the two alternatives. Consequently
\[
 e=1e=\sum_B b_Be=\sum_{B\in S(e)}b_B.
\]
Multiplication of a proposed sum by $b_C$ gives $b_C$ when $C$ is in
its index set and zero otherwise. Since $b_C\ne0$, this recovers the
index set from the sum and proves uniqueness. An empty sum is zero,
whereas a nonempty sum has a nonzero product with any of its members.
\end{proof}

\begin{implementation}
\module{ModularRep.CentralIdempotentBlockExpansion} proves the expansion
and its uniqueness in
\lean{ModularRep.CompleteOrthogonalCentralIdempotents.sum_centralIdempotentSupport_eq}
and \lean{ModularRep.CompleteOrthogonalCentralIdempotents.eq_centralIdempotentSupport_of_sum_eq}.
For a given block decomposition,
\module{ModularRep.BlockIdempotentDecomposition} defines
\lean{moduleBlock} and the equivalence \lean{primitiveBlockEquiv}
between the indexing set and the primitive central idempotents.
\end{implementation}

\begin{example}[Using a decomposition]
Given a simple representation and a block decomposition, the unique
idempotent acting as the identity determines its block. An isomorphism of modules
preserves this index because it preserves each equation $b_Bv=v$.
This is proved in
\lean{ModularRep.BlockIdempotentDecomposition.moduleBlock_eq_of_linearEquiv}.
\end{example}

\section{Idempotents of central subgroups}
\label{lib:blocks:central-sectors}

Let $Z\leq Z(G)$ have order invertible in $k$. For a linear character
$\nu:Z\to k^\times$, define
\[
 s_\nu=\sum_{z\in Z}\nu(z)^{-1}z,
 \qquad e_\nu=|Z|^{-1}s_\nu\in kG.
\]
The inverse in the coefficient ensures that $e_\nu$ acts as the identity
on a module with scalar action $zv=\nu(z)v$.

\begin{proposition}[Idempotents associated with central characters]
\label{lib:blocks:character-idempotents}
Each $e_\nu$ is a central idempotent, and $e_\nu e_\mu=0$ for
$\nu\ne\mu$. If $k$ is algebraically closed, then
\[
 \sum_{\nu\in\operatorname{Hom}(Z,k^\times)}e_\nu=1.
\]
Thus the idempotents form a complete family of orthogonal central idempotents.
\end{proposition}
\begin{proof}
For a nontrivial linear character $\chi$ of a finite group, choose $a$
with $\chi(a)\ne1$. Translation by $a$ permutes the group, so
\[
 \chi(a)\sum_z\chi(z)=\sum_z\chi(az)=\sum_z\chi(z).
\]
Cancellation in the field gives $\sum_z\chi(z)=0$. For the trivial
character the sum is the group order. Reindexing the product of the
weighted sums by $t=xy$ gives
\[
 s_\nu s_\mu
 =\left(\sum_{x\in Z}\nu(x)^{-1}\mu(x)\right)s_\mu.
\]
This is $|Z|s_\nu$ when $\nu=\mu$ and zero otherwise. Division by
$|Z|^2$ proves idempotence and orthogonality. Each group element in
the sum is central in $G$, which proves centrality in $kG$.

For completeness, finite abelian duality applies because $k$ has enough
roots of unity for the exponent of $Z$. Its character group has order
$|Z|$ and separates the elements of $Z$. If $z\ne1$, evaluation at
$z^{-1}$ is therefore a nontrivial character of this dual group. The
same translation argument gives
\[
 \sum_\nu\nu(z)^{-1}=0\quad(z\ne1),
 \qquad \sum_\nu\nu(1)^{-1}=|Z|.
\]
These are exactly the coefficient equations for $\sum_\nu e_\nu=1$.
Algebraic closure and invertibility of $|Z|$ provide the required roots.
\end{proof}

\begin{implementation}
The weighted product is
\lean{ModularRep.linearCharacterWeightedSum_mul} in
\module{ModularRep.LinearCharacterOrthogonality}.
The completeness proof is
\lean{ModularRep.sum_centralCharacterIdempotent_eq_one} in
\module{ModularRep.LinearCharacterCompleteness}. It uses Mathlib's finite
abelian duality to count the characters and separate group elements,
with the required roots of unity specified. The algebraically
closed version is proved in
\lean{ModularRep.completeCentralCharacterIdempotentsOfIsAlgClosed}.
\end{implementation}

For a finite dimensional irreducible representation $\rho$ over an
algebraically closed field, Schur's lemma gives a unique linear character
$\nu_\rho:Z\to k^\times$ such that
$\rho(z)=\nu_\rho(z)\id$. The block central character $\omega_B$
introduced below has domain $Z(kG)$, whereas $\nu_\rho$ has domain $Z$.

\begin{proposition}[The central character of a block]
\label{lib:blocks:sector}
Assume $k$ is algebraically closed and $|Z|$ is invertible in $k$.
Every primitive central idempotent $b$ lies above a unique character $\nu_b$
such that
\[
 be_{\nu_b}=b,\qquad be_\mu=0\quad(\mu\ne\nu_b).
\]
If $\rho$ is finite dimensional and irreducible, and $b$ acts as the
identity on $\rho$, then $\nu_\rho=\nu_b$. Hence any two such
representations have the same scalar action of $Z$.
\end{proposition}
\begin{proof}
Apply Proposition~\ref{lib:blocks:unique-support} to the complete
family of Proposition~\ref{lib:blocks:character-idempotents}.
On $\rho$, the action of $e_\mu$ is multiplication by
\[
 |Z|^{-1}\sum_{z\in Z}\mu(z)^{-1}\nu_\rho(z),
\]
which is one for $\mu=\nu_\rho$ and zero otherwise. Since $b$ acts as
the identity on $\rho$, Proposition~\ref{lib:blocks:unique-support}
gives $\nu_b=\nu_\rho$.
\end{proof}

\begin{implementation}
\lean{ModularRep.IsPrimitiveCentralIdempotent.existsUnique_centralCharacterSector}
and \lean{Representation.primitiveCentralIdempotentSector_eq_centralCharacter}
are in \module{ModularRep.CentralCharacterBlockSector}.
The scalar calculation is in \module{ModularRep.CentralCharacterAction}.
The equality of central characters assumes an irreducible representation
belonging to the block.
\end{implementation}

\begin{example}[Automorphisms and central characters]
Assume $k$ is algebraically closed. A ring isomorphism preserves
primitivity: pull back a central idempotent
below the image of $b$, apply primitivity of $b$, then map forward.
For a group automorphism $\alpha$, transport the basis of the group algebra by
$\alpha^{-1}$. If the order of $Z(G)$ is invertible in $k$, coefficient
comparison gives
\[
 e_\nu\longmapsto e_{\nu\circ\alpha|_{Z(G)}}.
\]
Thus a block lying above $\nu$ is sent to a block lying above
$\nu\circ\alpha|_{Z(G)}$. The corresponding theorem is
\lean{ModularRep.IsPrimitiveCentralIdempotent.centralCharacterSector_mapDomainRingEquiv_center}.
This example concerns the whole centre under the stated invertibility
hypothesis, as in \module{ModularRep.PrimitiveBlockAutomorphism}.
\end{example}

\section{Block central characters and induction}
\label{lib:blocks:induction}

Assume in this section that $k$ is algebraically closed, and fix a block
decomposition. The following standard description of block central
characters is used as an explicit input to the Lean argument.
In positive characteristic it is given by
\cite[Theorem~3.11]{Navarro1998}.

\begin{hypothesis}[Central characters of blocks]
\label{lib:blocks:catalogue}
For each $B\in\mathcal B(G)$ there is a $k$-algebra homomorphism
$\omega_B:Z(kG)\to k$ satisfying
\[
 \omega_B(b_C)=
 \begin{cases}1& B=C,\\0& B\ne C.\end{cases}
\]
Every $k$-algebra homomorphism $Z(kG)\to k$ is one of the $\omega_B$.
\end{hypothesis}

The evaluation equations imply that distinct blocks have distinct
central characters: evaluate an equality $\omega_B=\omega_C$ at $b_B$.
Thus $B\mapsto\omega_B$ is a bijection between the blocks and the
$k$-algebra homomorphisms $Z(kG)\to k$. The central character $\omega_B$
does not generally describe multiplication on the block algebra:
$zb_B=\omega_B(z)b_B$ need not hold for $z\in Z(kG)$.
\begin{implementation}
The hypotheses are expressed by
\lean{ModularRep.BlockCentralCharacterCatalogue}. The bijection is
\lean{ModularRep.BlockCentralCharacterCatalogue.centralCharacterEquiv}.
\end{implementation}

For $H\leq G$, coefficient restriction is the linear map
\[
 r_H:kG\longrightarrow kH,
 \qquad r_H\left(\sum_{g\in G}z_gg\right)=\sum_{h\in H}z_hh.
\]
It carries $Z(kG)$ into $Z(kH)$: the coefficients of a central
element are constant on $G$-conjugacy classes, so the restricted
coefficients are constant on $H$-conjugacy classes. Thus $r_H(z)$
commutes with every basis element of $kH$. Also $r_H(1)=1$.
Multiplicativity does not hold for a general subgroup. Even for an
abelian group and $H=1$, an element $g\ne1$ has
$r_H(g)=r_H(g^{-1})=0$ but $r_H(gg^{-1})=1$.

\begin{definition}
We use the definition of block induction by central characters
(see, for example, \cite[p.~87]{Navarro1998}). Given the block central characters for
$H$ and $G$, block induction from a block
$b\in\mathcal B(H)$ is \emph{defined} when the linear function
$\omega_b\circ r_H:Z(kG)\to k$ is multiplicative.
A block $B\in\mathcal B(G)$ is \emph{induced from} $b$ if this
condition holds and $\omega_B=\omega_b\circ r_H$.
\end{definition}

\begin{proposition}[Existence and uniqueness of an induced block]
\label{lib:blocks:induced-block}
If block induction from $b$ is defined, there is exactly one induced
block $B$.
\end{proposition}
\begin{proof}
The composite is $k$-linear, multiplicative and sends one to one. It is
therefore a $k$-algebra homomorphism. By
Hypothesis~\ref{lib:blocks:catalogue}, it equals $\omega_B$ for some $B$.
Uniqueness follows from injectivity of
$B\mapsto\omega_B$, proved by evaluation at block idempotents.
\end{proof}

\begin{implementation}
\module{ModularRep.GroupAlgebraCentralFunctions} uses Mathlib's group algebra
and its centre. It proves centrality of coefficient restriction and
constructs \lean{centerCoeffRestrict}.
\module{ModularRep.BlockInduction} defines
\lean{IsBlockInductionDefined}, proves
\lean{existsUnique_blockInducesTo}, and defines \lean{inducedBlock}.
The induced block is defined from multiplicativity and the given block
central characters.
The external source interfaces in
\module{ModularRep.SpathNavarroSourceAdapter} use an algebraically closed
field $k$ of characteristic $\ell$ that is algebraic over $\F_\ell$.
This additional hypothesis is recorded by
\lean{ModularRep.SpathCoefficientField}.
\end{implementation}

\begin{proposition}[Preservation on a central subgroup]
\label{lib:blocks:induction-sector}
Let $Z\leq Z(G)\cap H$, with $|Z|$ invertible in $k$.
Suppose $b\in\mathcal B(H)$ induces to $B\in\mathcal B(G)$.
Let $V$ and $W$ be finite dimensional irreducible representations of
$H$ and $G$, belonging to the blocks $b$ and $B$, respectively.
After identifying $Z$ with its copy in $H$, their scalar characters
on $Z$ agree.
\end{proposition}
\begin{proof}
For a central idempotent $e\in kG$, primitivity gives
$b_Be\in\{0,b_B\}$. Applying $\omega_B$ shows
\[
 \omega_B(e)=1\quad\Longleftrightarrow\quad b_Be=b_B.
\]
The reverse implication uses $\omega_B(b_B)=1$. In the forward
implication the alternative $b_Be=0$ would give $1=0$.

Choose $\nu$ so that $b_B$ lies above $\nu$ on $Z$. Since $Z\leq H$,
$e_\nu$ is supported on $H$, so $r_H(e_\nu)$ is the corresponding
character idempotent in
$kH$. In particular it is still idempotent. Induction gives
\[
 \omega_b(r_H(e_\nu))=\omega_B(e_\nu)=1.
\]
The same criterion in $kH$ gives
$b_b r_H(e_\nu)=b_b$. Hence $b_b$ also lies above $\nu$.
Proposition~\ref{lib:blocks:sector} identifies $\nu$ with
the scalar characters of $V$ and $W$.
\end{proof}

\begin{implementation}
The criterion for the value on a central idempotent is
\lean{ModularRep.BlockCentralCharacterCatalogue.centralIdempotent_value_one_iff}.
The conclusion is
\lean{Representation.centralCharacter_eq_comp_of_blockInducesTo}
in \module{ModularRep.CentralCharacterBlockInductionRestriction}.
The proof identifies $Z$ with its copy in $H$ and uses the resulting
equality of orders. It does not require $V$ to occur in the restriction
of $W$.
\end{implementation}

\section{The central Brauer map}
\label{lib:blocks:brauer-map}

Now assume $\operatorname{char}k=\ell$ and let $P\leq G$ be an
$\ell$-subgroup. Write $C=C_G(P)$. The \emph{central Brauer map}
is coefficient restriction to $C$ on the centre of $kG$:
\[
 \operatorname{Br}_P:Z(kG)\longrightarrow Z(kC),
 \qquad \operatorname{Br}_P(z)=\sum_{c\in C}z_cc.
\]
The target is the centre of $kC$, by the restriction argument above.
Its multiplicativity follows from an orbit calculation. For this case
$H=C_G(P)$, see, for example, \cite[Theorem~4.9]{Navarro1998}.

\begin{lemma}[Summing over fixed points]
\label{lib:blocks:orbit-sum}
Let an $\ell$-group $Q$ act on a finite set $X$. If $f:X\to k$ is
constant on each $Q$-orbit, then
\[
 \sum_{x\in X}f(x)=\sum_{x\in X^Q}f(x).
\]
\end{lemma}
\begin{proof}
Partition $X$ into orbits. On an orbit $\mathcal O$, the sum equals
$|\mathcal O|f(x)$ for any $x\in\mathcal O$. Orbit sizes are powers
of $\ell$, so an orbit of size greater than one contributes zero in
characteristic $\ell$. An orbit of size one is exactly a fixed point
and contributes its own value.
\end{proof}

\begin{theorem}[The Brauer map on the centre]
\label{lib:blocks:brauer-multiplicative}
For a finite group $G$, an $\ell$-subgroup $P$ and a field $k$ of
characteristic $\ell$, the map $\operatorname{Br}_P$ is a
$k$-algebra homomorphism. Algebraic closure is not required.
\end{theorem}
\begin{proof}
Fix $c\in C$. The coefficient of $zw$ at $c$ is
\[
 (zw)_c=\sum_{a\in G}z_a w_{a^{-1}c}.
\]
Let $P$ act on $G$ by conjugation. If $q\in P$, then $qc=cq$, so
\[
 (qaq^{-1})^{-1}c=q(a^{-1}c)q^{-1}.
\]
Centrality of $z$ and $w$ makes the function
$a\mapsto z_aw_{a^{-1}c}$ constant on these orbits. Its fixed points
are precisely $a\in C$. Lemma~\ref{lib:blocks:orbit-sum} gives
\[
 (zw)_c=\sum_{a\in C}z_aw_{a^{-1}c}
       =\bigl(\operatorname{Br}_P(z)\operatorname{Br}_P(w)\bigr)_c.
\]
Here $a^{-1}c\in C$ whenever $a,c\in C$, so the final expression is
the product coefficient in $kC$. Equality at every $c\in C$ proves
multiplicativity. Linearity and preservation of one were already proved
for coefficient restriction.
\end{proof}

\begin{implementation}
The proof uses Mathlib's theorem that orbit sizes for an $\ell$-group
are powers of $\ell$. The resulting sum formula is
\lean{ModularRep.IsPGroup.sum_eq_sum_fixedPoints_of_invariant}
in \module{ModularRep.PGroupInvariantSum}.
Multiplicativity and the algebra homomorphism are
\lean{ModularRep.centralBrauerRestriction_mul} and
\lean{ModularRep.centralBrauerMap} in
\module{ModularRep.GroupAlgebraCentralBrauerMap}.
These declarations allow any commutative coefficient semiring of
characteristic $\ell$.
\end{implementation}

\begin{proposition}[The Brauer map to an intermediate subgroup]
\label{lib:blocks:interval-map}
Suppose $C_G(P)\leq H\leq N_G(P)$. Extend the coefficients of
$\operatorname{Br}_P(z)$ by zero from $C_G(P)$ to $H$.
The resulting map
\[
 \operatorname{Br}_P^H:Z(kG)\longrightarrow Z(kH)
\]
is a $k$-algebra homomorphism. Moreover,
$\operatorname{Br}_P(z)\ne0$ if and only if
$\operatorname{Br}_P^H(z)\ne0$.
\end{proposition}
\begin{proof}
The inclusion $kC_G(P)\hookrightarrow kH$ is an injective algebra
homomorphism. Every $h\in H$ normalises $P$, and therefore preserves
$C_G(P)$ by conjugation. The coefficients $z_c$ are constant
under this conjugation because $z\in Z(kG)$. The extended element
thus lies in $Z(kH)$. Multiplication and the unit are preserved by the
composite. Injectivity gives the asserted equivalence of nonvanishing.
\end{proof}

For this form of the Brauer homomorphism, see, for example,
\cite[Theorem~4.9]{Navarro1998}. Its image is central in $kH$
because the restricted coefficients are invariant under $H$.
An arbitrary element of $Z(kC_G(P))$ need not have this property.

\begin{implementation}
The homomorphism \lean{ModularRep.centralBrauerMapTo} is in
\module{ModularRep.GroupAlgebraCentralBrauerMapTo}.
The equivalence of nonvanishing,
\lean{ModularRep.centralBrauerMap_ne_zero_iff_mapTo_ne_zero},
is in \module{ModularRep.CentralBrauerIntervalSupport}.
For $H=N_G(P)$, the map is named
\lean{ModularRep.normalizerCentralBrauerMap}.
\end{implementation}

\begin{example}[A local sum of blocks]
The image of a primitive central idempotent of $kG$ under
$\operatorname{Br}_P^H$ is a central idempotent, although it can be
zero or a sum of several primitive central idempotents of $kH$. A block
decomposition of $kH$ and
Proposition~\ref{lib:blocks:idempotent-expansion} determine its unique
support. This is implemented by
\lean{ModularRep.sum_centralBrauerMapToLocalSupport_eq} and
\lean{ModularRep.centralBrauerMapToPrimitiveImage_ne_zero_iff_support_nonempty}.
In particular, the Brauer homomorphism need not preserve primitivity.
\end{example}

\section{Brauer's first main theorem}
\label{lib:blocks:normaliser-correspondence}

Brauer's first main theorem identifies the blocks of $G$ with defect
group $P$ with the blocks of $N_G(P)$ having defect group $P$.
See, for example, \cite[Theorem~4.17]{Navarro1998}.
We prove the correspondence from the central character identity and the
five facts about defect groups stated below.
Assume $k$ is algebraically closed of
characteristic $\ell$. Fix an $\ell$-subgroup $P$, write
$N=N_G(P)$, and fix block decompositions and block central characters
for $N$ and $G$ as in Hypothesis~\ref{lib:blocks:catalogue}.

We first use the central character identity
(see, for example, \cite[Theorem~4.14]{Navarro1998}). For
$PC_G(P)\leq H\leq N_G(P)$, it is
\begin{equation}
 \omega_b\circ r_H=\omega_b\circ\operatorname{Br}_P^H
 \qquad(b\in\mathcal B(H)).
 \label{lib:blocks:414-equality}
\end{equation}
The hypothesis includes $P\leq H$, whereas the construction of
$\operatorname{Br}_P^H$ only required $C_G(P)\leq H\leq N_G(P)$.
Since the right side of \eqref{lib:blocks:414-equality} is an algebra
homomorphism, block induction is defined. For $H=N$, denote the induced
block map by
$I:\mathcal B(N)\to\mathcal B(G)$. Thus
\begin{equation}
 \omega_{I(b)}=\omega_b\circ\operatorname{Br}_P^N.
 \label{lib:blocks:induction-brauer-character}
\end{equation}

For the following identity, see, for example, the final assertion of
\cite[Theorem~4.14]{Navarro1998}, with $H=N_G(P)$.

\begin{proposition}[The Brauer image as a sum of blocks]
\label{lib:blocks:fibre}
Under \eqref{lib:blocks:414-equality}, for every block $B$ of $G$,
\[
 \operatorname{Br}_P^N(b_B)=\sum_{b\in\mathcal B(N):\ I(b)=B}b_b.
\]
In particular, if the image is nonzero, some block of $N$ induces to $B$.
\end{proposition}
\begin{proof}
Expand the image as in Proposition~\ref{lib:blocks:idempotent-expansion}.
A local block $b$ lies in its support exactly when its central
character takes value one on the image. Indeed, multiplication by $b_b$
is either zero or $b_b$, and evaluation by $\omega_b$ distinguishes
these alternatives. Equation~\eqref{lib:blocks:induction-brauer-character}
identifies this value with $\omega_{I(b)}(b_B)$. The evaluation equations
for the ambient block central characters say that this is one exactly
when $I(b)=B$.
\end{proof}

\begin{implementation}
\module{ModularRep.Navarro414IntervalCentralCharacterSource} expresses
\eqref{lib:blocks:414-equality}, from which the induced block is defined.
The equivalence between occurrence in this sum and induction to $B$ is
\lean{ModularRep.mem_navarro417PhaseANormalizerSupport_iff_inductionMap_eq}
in \module{ModularRep.Navarro417PhaseASupport}.
\end{implementation}

For the Brauer map characterisation of defect groups, see, for example,
\cite[Theorem~4.11]{Navarro1998}. For a block $B$ of a finite group
$X$, let $\mathcal D_X(B)$ be the set of $\ell$-subgroups $D\leq X$
such that
\begin{equation}
 \operatorname{Br}_Q(b_B)\ne0
 \quad\Longleftrightarrow\quad
 Q\text{ is conjugate to a subgroup of }D
 \quad\text{for every }\ell\text{-subgroup }Q\leq X.
 \label{lib:blocks:defect-support}
\end{equation}
In the following argument, we use this criterion to define
$\mathcal D_X(B)$. Its members are the defect groups of $B$.
Any two members are conjugate: the criterion makes each conjugate to
a subgroup of the other, and equality of their finite orders then gives
conjugacy. Existence will be used explicitly below.

\begin{hypothesis}[Standard facts for the defect correspondence]
\label{lib:blocks:defect-hypotheses}
We use the following five standard facts about the fixed block
decompositions and central characters.
\begin{enumerate}
\item If $P\in\mathcal D_N(b)$ and an $N$-conjugacy class
$\mathcal K$ has nonzero coefficient in $b_b$, then
$\mathcal K=\mathcal L\cap C_G(P)$ for some $G$-conjugacy class
$\mathcal L$, with the intersection viewed inside $N$.
\item The set $\mathcal D_N(b)$ is nonempty for every block $b$ of $N$.
\item If $D\in\mathcal D_N(b)$, then $O_\ell(N)\leq D$.
\item If $b$ induces to $B$ and $D\in\mathcal D_N(b)$, then there is
$E\in\mathcal D_G(B)$ containing $D$, viewed as a subgroup of $G$.
\item If $P\in\mathcal D_N(b)$, then there is
$E\in\mathcal D_G(I(b))$ with $E\leq P$.
\end{enumerate}
\end{hypothesis}

For the first condition, Lemma~4.15 restricts the nonzero class
coefficients to classes with defect group $P$, and Lemma~4.16 writes
each such class as the required intersection
\cite[pp.~88--89]{Navarro1998}. The second condition follows from
existence of defect groups and their Brauer map characterisation
\cite[Theorems~4.3 and~4.11]{Navarro1998}.
The third and fourth conditions follow from
\cite[Theorem~4.8 and Lemma~4.13]{Navarro1998},
respectively.

The fifth condition is the containment proved in the first paragraph
of the proof of \cite[Theorem~4.17]{Navarro1998}. That argument
chooses a defect class of $b$, uses Corollary~4.5 and Lemma~4.16 to
lift it to an ambient class with defect group $P$, and applies
Theorems~4.14 and~4.4 to obtain a defect group of $I(b)$ contained in
$P$. We use this containment, together with the other four conditions
and \eqref{lib:blocks:414-equality}, to prove the correspondence.

\begin{proposition}[Lifting a local block idempotent]
\label{lib:blocks:class-lift}
Assume the first condition of
Hypothesis~\ref{lib:blocks:defect-hypotheses} and
\eqref{lib:blocks:414-equality}. If $P\in\mathcal D_N(b)$, then
$b_b$ is in the image of $\operatorname{Br}_P^N$.
If $I(b)=I(c)$ for any local block $c$, then $b=c$. Consequently,
\[
 \operatorname{Br}_P^N(b_{I(b)})=b_b.
\]
\end{proposition}
\begin{proof}
For an $N$-class $\mathcal K$, write
$\mathcal K^+=\sum_{x\in\mathcal K}x$. Centrality gives a unique
class expansion
\[
 b_b=\sum_{\mathcal K}a_{\mathcal K}\mathcal K^+.
\]
For each nonzero coefficient choose an ambient class
$\mathcal L_{\mathcal K}$ with
$\mathcal L_{\mathcal K}\cap C_G(P)=\mathcal K$.
Coefficient restriction gives
$\operatorname{Br}_P^N(\mathcal L_{\mathcal K}^+)=\mathcal K^+$.
Thus the finite sum
\[
 z=\sum_{a_{\mathcal K}\ne0}
       a_{\mathcal K}\mathcal L_{\mathcal K}^+\in Z(kG)
 \quad\text{satisfies}\quad \operatorname{Br}_P^N(z)=b_b.
\]

If $I(b)=I(c)$, equation~\eqref{lib:blocks:induction-brauer-character}
gives equal composites
$\omega_b\operatorname{Br}_P^N=\omega_c\operatorname{Br}_P^N$.
Evaluation at $z$ gives $\omega_b(b_b)=\omega_c(b_b)$. The left side
is one, and the right side is zero unless $b=c$. This proves
injectivity even when only the first block is known to have defect
$P$. In the sum of Proposition~\ref{lib:blocks:fibre}, the only block
inducing to $I(b)$ is therefore $b$. This gives the final equation.
\end{proof}

\begin{implementation}
The element $z$ is constructed in
\module{ModularRep.Navarro415416ClassSumAdapter}, in the theorem
\lean{ModularRep.Navarro415416RawClassSource.localBlockIdempotent_mem_normalizerCentralBrauerMap_range}.
Injectivity and the equality of block idempotents are
\lean{ModularRep.navarro417PhaseAInductionMap_eq_imp_eq_of_localP_left}
and \lean{ModularRep.normalizerCentralBrauerMap_selectedBlockIdempotent_eq}
in \module{ModularRep.Navarro417RestrictedInjectivity}.
The class intersection property and the central character equality
suffice for all three conclusions.
\end{implementation}

\begin{theorem}[Blocks with a prescribed defect group]
\label{lib:blocks:exact-defect}
Assume \eqref{lib:blocks:414-equality} and
Hypothesis~\ref{lib:blocks:defect-hypotheses}. Write
$\mathcal B(X\mid P)=\{B\in\mathcal B(X):P\in\mathcal D_X(B)\}$,
where $X=G$ or $N_G(P)$ and $P$ is viewed as a subgroup of $X$.
Then block induction is a bijection
\[
 I:\mathcal B(N_G(P)\mid P)\longrightarrow\mathcal B(G\mid P).
\]
\end{theorem}
\begin{proof}
Let $b\in\mathcal B(N\mid P)$. The fourth condition gives
$E_+\in\mathcal D_G(I(b))$ with $P\leq E_+$.
The fifth gives $E_-\in\mathcal D_G(I(b))$ with $E_-\leq P$.
They are conjugate by
\eqref{lib:blocks:defect-support}, so
\[
 |P|\leq|E_+|=|E_-|\leq|P|.
\]
The inclusions are therefore equalities. In particular,
$P\in\mathcal D_G(I(b))$.

Conversely, suppose $I(b)=B$ and $P\in\mathcal D_G(B)$.
By the second condition, choose $D\in\mathcal D_N(b)$. Since
$P$ is a normal $\ell$-subgroup of $N$, the third condition gives
\[
 P\leq O_\ell(N)\leq D.
\]
The fourth condition gives $E\in\mathcal D_G(B)$ with
$D\leq E$ inside $G$. The defect groups $E$ and $P$ are conjugate. Hence
$|P|\leq|D|\leq|E|=|P|$, and the local inclusion is an equality.
Thus $P\in\mathcal D_N(b)$ for every local block inducing to $B$.

Finally, if $P\in\mathcal D_G(B)$, then
$\operatorname{Br}_P(b_B)\ne0$ by
\eqref{lib:blocks:defect-support}. Proposition~\ref{lib:blocks:interval-map}
preserves this nonvanishing in $kN$, and
Proposition~\ref{lib:blocks:fibre} supplies a local block over $B$.
The preceding paragraph puts that block in $\mathcal B(N\mid P)$,
proving surjectivity. Injectivity follows from
Proposition~\ref{lib:blocks:class-lift}.
\end{proof}

\begin{implementation}
\module{ModularRep.Navarro417ExactDefectAdapter} proves the two
containment arguments and the bijection as
\lean{navarro417PhaseAInductionMap_has_ambientP},
\lean{navarro417LocalP_of_blockInducesTo_ambientP}, and
\lean{navarro417ExactDefectMap_bijective}, all in namespace
\lean{ModularRep}. The resulting equivalence is
\lean{ModularRep.navarro417ExactDefectEquiv}.
The construction takes the five conditions above and
\eqref{lib:blocks:414-equality}, for the given block decompositions
and central characters. The third condition, containment of the normal
$\ell$-core in every defect group, is proved by
\lean{navarro408PCoreDefectSource} in
\module{ModularRep.PaperProofs.SporadicFi24P3Definition44NamedCarrierPCoreDefectContainment}.
The remaining source
conditions used by this construction are stated separately.
\end{implementation}

\section{Linear characters stabilising a block}
\label{lib:blocks:linear-stabilisers}

Let $N\trianglelefteq G$ and assume again that $k$ is algebraically
closed of characteristic $\ell$. Write $\pi:G\to G/NZ(G)$ for the
quotient homomorphism. A linear character
$\lambda:G\to k^\times$ twists a representation by
\[
 (\lambda\otimes\rho)(g)=\lambda(g)\rho(g).
\]
If a central subgroup $Z$ acts on $\rho$ by a scalar character
$\nu_\rho$, it acts on $\lambda\otimes\rho$ by
$\lambda|_Z\nu_\rho$. This follows from the scalar action and does
not require division by the degree of $\rho$.

Suppose a group $A$ acts on the primitive central idempotents of $kG$, and there
is an injective homomorphism
$a\mapsto\lambda_a$ from $A$ to the linear characters of $G$.
Assume every $\lambda_a$ is trivial on $N$. Suppose every block contains
a finite dimensional simple representation, and tensoring any such
representation in $b$ by $\lambda_a$ gives a representation in
$a\cdot b$.

\begin{theorem}[Stabilisers under tensoring by linear characters]
\label{lib:blocks:stabiliser-bound}
Under these assumptions, the stabiliser
$A_b=\{a\in A:a\cdot b=b\}$ is finite. For every $a\in A_b$, the
character $\lambda_a$ factors through $G/NZ(G)$, and
\[
 |A_b|\leq [G:NZ(G)].
\]
\end{theorem}
\begin{proof}
The central $\ell$-regular elements form a subgroup
$Z_{\ell'}\leq Z(G)$ of order prime to $\ell$.
Choose a simple representation $\rho$ belonging to $b$. If $a\in A_b$,
then $\lambda_a\otimes\rho$ belongs to the same block.
Proposition~\ref{lib:blocks:sector} gives equality of their scalar characters on
$Z_{\ell'}$. Since twisting multiplies that character by
$\lambda_a|_{Z_{\ell'}}$, cancellation in $k^\times$ proves that
$\lambda_a$ is trivial on $Z_{\ell'}$.

It is also trivial on every $\ell$-element. If $g^{\ell^m}=1$, then
$\lambda_a(g)^{\ell^m}=1$. In characteristic $\ell$,
$X^{\ell^m}-1=(X-1)^{\ell^m}$, so a field has no nontrivial
$\ell$-power roots of unity. Every central element is the product
of its commuting $\ell$-part and $\ell$-regular part, both powers
of the original element. Hence $\lambda_a$ is trivial on all of
$Z(G)$, and therefore on $NZ(G)$.

Thus $\lambda_a$ induces a character
$\overline\lambda_a:G/NZ(G)\to k^\times$ with
$\lambda_a=\overline\lambda_a\circ\pi$. These quotient characters
are distinct, since composing with $\pi$ recovers $\lambda_a$ and
$a\mapsto\lambda_a$ is injective. These quotient characters form a
linearly independent family in the space of all functions on
$G/NZ(G)$, which has dimension $[G:NZ(G)]$. Therefore their indexing
set $A_b$ is finite and $|A_b|\leq [G:NZ(G)]$.
\end{proof}

\begin{implementation}
The assumptions on simple representations and tensoring are expressed by
\lean{ModularRep.LinearBlockStabilizer.SimpleBlockTensorSource}.
It supplies a simple representation in every block and compatibility
between the block action and tensoring. The theorem is
\lean{ModularRep.LinearBlockStabilizer.modular_block_stabilizer_bound}.
The quotient character is defined by \lean{descendCenterQuotient}.
Mathlib's linear independence of multiplicative characters gives
\lean{finite_of_injective_linear_characters} and the bound
\lean{card_le_of_injective_linear_characters}, in the same namespace.
The theorem \lean{block_stabilizer_finite} explicitly supplies finiteness
of the stabiliser for both the ordinary and modular character actions.
\end{implementation}

\begin{example}[All quotient characters]
Take $A$ to be the linear characters of $G$ trivial on $N$, identified
with those of $G/N$. Once their block action and its compatibility with
tensoring are given, every character stabilising a block factors through
$G/NZ(G)$, and the stated bound follows. In the special case $NZ(G)=G$,
the bound is one, so every block has trivial stabiliser for this
action.
\end{example}

For the ordinary form, let $K$ be a field of characteristic zero.
Set $m=|G|_{\ell'}$ and fix an isomorphism
$\mu_m(k)\simeq\mu_m(K)$ between the groups of $m$th roots of unity.
Applied pointwise, it identifies ordinary linear characters of $G/N$
of order prime to $\ell$ with modular linear characters of $G/N$.
Transport the block action through this identification. The stabilisers
are isomorphic, and preservation of kernels gives the factorisation
through $G/NZ(G)$ and the same bound.

\begin{implementation}
\module{ModularRep.LinearBlockStabilizerReduction} constructs
\lean{ordinaryToBrauerQuotient}, \lean{reductionStabilizerEquiv}, and
\lean{ordinary_block_stabilizer_bound} in namespace
\lean{ModularRep.LinearBlockStabilizer}.
The theorem \lean{reduction_commutes_with_tensor} requires a modular system,
\lean{StableReductionBrauerCharacterCompatibility} and
\lean{BrauerLinearTensorProductFormula}.
\end{implementation}
